\scanpage{373}
\mspage{112}
% archive sheet number: 372
ou de $\mathrm{Ind\,Hot}$. (Dont la limite inductive dans $\mathrm{Hot}_0^{\unsure{\wedge|1}}$ \unsure{n'est} autre que l'élément de $\mathrm{Hot}_0^\wedge$ défini par $C_\infty$.)

Mes efforts pour prouver que \deleted{\unsure{$C$}}
\[
C_\infty\mathrel{\unsure{-|\to}}C
\]
était un iso. dans $\mathrm{Hot}$, ou ne serait-ce qu'un $1$-iso (i.e. un isom. non seulement pour les $\pi_0$, ce qui est tjrs vrai, mais aussi pour les $\pi_1$ de \deleted{\unsure{toutes}} \unsure{les|ces} composantes), ont échoué. J'ai trouvé qu'il \unsure{suffirait|suffisait} que le foncteur
\[
T\longrightarrow T_0
\]
soit ou bien lisse et coasphérique, ou bien propre et asphérique -- donc \unsure{tout est} OK si $T=T_0$ i.e. si $T$ est ordonné. Mais dans le cas $T=\unsure{\underline{\mathrm{Tg}}|\underline{T_g}}$, le foncteur $T\mathrel{-}T_0$ n'est ni propre, ni lisse (\unsure{pas|pr} $=$ $\unsure{\mathbb D|D|0}$-propre ou $\mathbb D$-lisse).

Finalement j'en suis arrivé à la conviction que dans le cas qui nous occupe $(\deleted{\unsure{C}}T=\underline{\mathrm{Tg}},\ C=\underline{\mathrm{Ch}}_{\unsure{\Sigma_{\mathscr B}|\Sigma_\beta}}(X;a,b))$ la \deleted{foncteur} \inserted{flèche} $C_\infty\longrightarrow C$ n'est en général \uline{pas} un \unsure{iso} \unsure{dans} $\mathrm{Hot}$. En effet, dans le cas $a=b$, \deleted{cela} si cela
